% !TeX root = ../main.tex
%%% ---------------------------------------------------------------------------
%%% 成果与工作量
%%%
%%% 目标页数：1 页。
%%% 这一节是学位答辩独有的，会议报告绝对不会有。委员会需要它来核对
%%% 学位授予的硬性条件（论文数、录用状态、专利、软著）。
%%%
%%% 写法要点：
%%%   · 论文按「作者顺序、刊物/会议全称、级别、状态」列，不要只写缩写
%%%   · 状态要写实：已发表 / 已录用 / 在审 / 已投稿，不要模糊成「发表」
%%%   · 代码、数据集、软著也算工作量，有就列
%%%   · 如果某项还在审，说清楚投稿时间，委员会会据此判断进度合理性
%%% ---------------------------------------------------------------------------
\section{成果与工作量}

\begin{frame}{攻读学位期间的研究成果}
  \begin{block}{学术论文}
    \begin{enumerate}
      \setlength{\itemsep}{0.15em}
      \item \textbf{张三}, 李四. 基于多尺度特征融合的遥感图像目标检测方法[J].
            自动化学报, 2026, 52(4): 801--815.（一作，中科院一区，\alert{已发表}）
      \item \textbf{Zhang S}, Li S. Deformable Alignment for Small Object
            Detection[C]. CVPR 2026.（一作，CCF-A，\alert{已录用}）
      \item Wang W, \textbf{Zhang S}. Survey on Remote Sensing Detection[J].
            IEEE TGRS.（二作，\alert{在审，2026 年 3 月投稿}）
    \end{enumerate}
  \end{block}

  \begin{block}{其他}
    \begin{itemize}
      \item 发明专利 1 项（已受理，申请号 CN2026XXXXXXX）
      \item 开源代码 \url{https://github.com/example/msff-net}，
            \num{420} stars
    \end{itemize}
  \end{block}

  \note{念这一页要慢，尤其是级别和状态。委员会在这一页做记录。}
\end{frame}
